\documentclass[a4paper]{article}

% --- Gói ngôn ngữ và font ---
\usepackage{vntex}             % Hỗ trợ tiếng Việt
\usepackage[english]{babel}    % Hỗ trợ tiếng Anh (tùy chọn)
\usepackage[utf8]{inputenc}
\usepackage{amsmath}
\usepackage{amssymb}
\usepackage{latexsym}
\usepackage{indentfirst} % Ép buộc thụt đầu dòng đoạn đầu tiên

% --- Gói hình ảnh và màu sắc ---
\usepackage{graphicx}
\usepackage{color}
\usepackage{xcolor} % Dùng xcolor thay cho color để mạnh mẽ hơn
\usepackage{svg}
\usepackage{tikz}
\usetikzlibrary{arrows,snakes,backgrounds}
\usepackage[framemethod=tikz]{mdframed}

% --- Gói định dạng trang và bảng ---
\usepackage{geometry}
\usepackage{a4wide}
\usepackage{fancyhdr}
\usepackage{booktabs}
\usepackage{array}
\usepackage{hhline}
\usepackage{multicol}
\usepackage{multirow}
\usepackage{makecell}
\usepackage{lastpage}
\usepackage{placeins}
\usepackage{setspace}
\usepackage[skip=0pt]{caption}
\usepackage{float} % Để cố định hình ảnh với [H]

% --- Gói code và thuật toán ---
\usepackage{algorithm}
\usepackage{algpseudocode}
\usepackage{listings}
\usepackage[most]{tcolorbox}
\usepackage{minted} % Yêu cầu -shell-escape

% --- Các gói khác ---
\usepackage{blindtext}
\usepackage{comment}
\usepackage{enumerate}
\usepackage{enumitem}
\usepackage{hyperref}

% --- Cấu hình link ---
\hypersetup{
    urlcolor=blue,
    linkcolor=black,
    citecolor=black,
    colorlinks=true
}

% --- Định nghĩa màu sắc ---
\setlength{\columnseprule}{1pt}
\def\columnseprulecolor{\color{black}}
\definecolor{codegreen}{rgb}{0,0.6,0}
\definecolor{codegray}{rgb}{0.5,0.5,0.5}
\definecolor{codepurple}{rgb}{0.58,0,0.82}
\definecolor{backcolour}{rgb}{0.95,0.95,0.92}
\definecolor{ReportNavy}{HTML}{102A43}
\definecolor{ReportCyan}{HTML}{32B5D2}
\definecolor{CodeBackground}{HTML}{F6F8FA}
\definecolor{CodeFrame}{HTML}{CBD5E1}
\definecolor{CodeComment}{HTML}{2E7D32}
\definecolor{CodeKeyword}{HTML}{005CC5}
\definecolor{CodeString}{HTML}{B31D28}
\definecolor{CodeNumber}{HTML}{64748B}

\newcommand{\exerciseOneURL}{https://github.com/USERNAME/REPOSITORY}
\newcommand{\exerciseTwoURL}{https://github.com/USERNAME/REPOSITORY}
\newcommand{\exerciseThreeURL}{https://github.com/USERNAME/REPOSITORY}

% --- Cấu hình Minted (C++) ---
\usemintedstyle{vs}
\setminted[cpp]{
    breaklines=true,
    mathescape=true,
    linenos=true,
    fontsize=\small,
    numbersep=3pt,
    bgcolor=pink!5,
    frame=single,
    framerule=0.5pt,
    rulecolor=pink!200
}

% --- Cấu hình Header/Footer ---
\everymath{\color{blue}}
\setlength{\headheight}{40pt}
\pagestyle{fancy}
\fancyhead{} % clear all header fields
\fancyhead[L]{
    \begin{tabular}{rl}
        \begin{picture}(25,15)(0,0)
            \put(0,-8){\includegraphics[width=8mm, height=8mm]{hcmut.png}}
        \end{picture}&
        \begin{tabular}{l}
            \textbf{\ttfamily University of Technology}\\
            \textbf{\ttfamily Faculty of Computer Science and Engineering}
        \end{tabular} 	
    \end{tabular}
}
\fancyhead[R]{} % Bên phải header để trống

\fancyfoot{} % clear all footer fields
\fancyfoot[L]{\scriptsize \ttfamily \textcolor{blue}{Discrete Structure (CO1007) Assignment, Semester 1, Academic year 2025-2026}}
\fancyfoot[R]{\scriptsize \ttfamily Page {\thepage}/\pageref{LastPage}}
\renewcommand{\headrulewidth}{0.3pt}
\renewcommand{\footrulewidth}{0.3pt}

% --- Cấu hình Mục lục ---
\setcounter{secnumdepth}{4}
\setcounter{tocdepth}{3}
\makeatletter
\newcounter {subsubsubsection}[subsubsection]
\renewcommand\thesubsubsubsection{\thesubsubsection .\@alph\c@subsubsubsection}
\newcommand\subsubsubsection{\@startsection{subsubsubsection}{4}{\z@}%
                                     {-3.25ex\@plus -1ex \@minus -.2ex}%
                                     {1.5ex \@plus .2ex}%
                                     {\normalfont\normalsize\bfseries}}
\newcommand*\l@subsubsubsection{\@dottedtocline{3}{10.0em}{4.1em}}
\newcommand*{\subsubsubsectionmark}[1]{}
\makeatother

\setlength{\parindent}{1cm}
\setlength{\parskip}{0pt}

\begin{document}

% --- TRANG BÌA ---
\begin{titlepage}
\begin{center}
    \Large Vietnam National University, Ho Chi Minh City \\
    \Large University of Technology \\
    \Large Faculty of Computer Science and Engineering 
\end{center}

\vspace{1cm}

\begin{figure}[H]
    \begin{center}
        \includegraphics[width=4cm]{hcmut.png}
    \end{center}
\end{figure}

\vspace{1cm}

\begin{center}
    \begin{tabular}{c}
    \multicolumn{1}{l}{\textbf{{\Large Microprocessors-Microcontrollers (Lab) (CO3010)}}}\\
    ~~\\
    \hline
    \\
    \multicolumn{1}{l}{\textbf{{\Large Assignment }}}\\
    \\
    \textbf{{\LARGE LAB\_01}}\\
    \\
    \hline
    \end{tabular}
\end{center}

\vspace{2cm}

\begin{table}[H]
    \centering
    \begin{tabular}{rrl}
        & Instructor(s): & Phan Văn Sỹ, CSE-HCMUT \\
        & Class: & CC03\\
        & Student: & Vo Le Thanh Trieu - 2453297\\
    \end{tabular}
\end{table}

\vspace{1cm}
\begin{center}
    {\footnotesize Ho Chi Minh City, September 2026}
\end{center}
\end{titlepage}

\thispagestyle{fancy}
\newpage
\tableofcontents
\newpage

%%%%%%%%%%%%%%%%%%%%%%%%%%%%%%%%%
% LƯU Ý: Hãy đổi tên file và thư mục của bạn thành KHÔNG DẤU KHÔNG CÁCH
% Ví dụ: Đề bài.tex -> DeBai.tex
% Dùng \input an toàn hơn \include cho các đoạn nhỏ
%%%%%%%%%%%%%%%%%%%%%%%%%%%%%%%%%

% ------------------------------------------------------------
% Astonishing STM32 C code block
% ------------------------------------------------------------
\newtcblisting[auto counter,number within=section]{stmcode}[1]{
    enhanced,
    breakable,
    listing only,
    colback=CodeBackground,
    colframe=CodeFrame,
    colbacktitle=ReportNavy,
    coltitle=white,
    fonttitle=\bfseries,
    title={Program~\thetcbcounter: #1},
    attach boxed title to top left={xshift=4mm,yshift=-2mm},
    boxed title style={
        colback=ReportNavy,
        colframe=ReportNavy,
        arc=1.5mm,
        boxrule=0pt
    },
    arc=2mm,
    boxrule=0.7pt,
    left=7mm,
    right=2mm,
    top=4mm,
    bottom=2mm,
    drop shadow southeast,
    before skip=1.2em,
    after skip=1.2em,
    listing options={
        language=C,
        basicstyle=\ttfamily\small,
        keywordstyle=\color{CodeKeyword}\bfseries,
        commentstyle=\color{CodeComment}\itshape,
        stringstyle=\color{CodeString},
        numbers=left,
        numberstyle=\tiny\color{CodeNumber},
        numbersep=10pt,
        stepnumber=1,
        showstringspaces=false,
        breaklines=true,
        breakatwhitespace=false,
        keepspaces=true,
        tabsize=4,
        columns=fullflexible,
        escapeinside={(*@}{@*)}
    }
}

% ------------------------------------------------------------
% Reusable figure placeholder
% The real image is loaded automatically when the file exists.
% ------------------------------------------------------------
\newcommand{\schematicplaceholder}[2]{%
    \IfFileExists{#1}{%
        \includegraphics[width=#2\textwidth]{#1}%
    }{%
        \fcolorbox{ReportCyan}{CodeBackground}{%
            \parbox[c][5.5cm][c]{0.82\textwidth}{%
                \centering
                \textcolor{ReportNavy}{\Large\bfseries Schematic Placeholder}\\[0.8em]
                Add the image file:\\[0.4em]
                \texttt{\detokenize{#1}}
            }%
        }%
    }%
}

% ============================================================
\section{LED Animation Exercises}

This report presents the design and implementation of the first three LED
animation exercises using the STM32F103C6Ux microcontroller. STM32CubeMX was
used to configure the GPIO pins and generate the initialization code. The
application code was edited and compiled in Visual Studio Code using CMake.

In all three exercises, the positive terminal of each LED is connected to the
3.3-V supply, while its negative terminal is connected to an STM32 GPIO output.
The LEDs are therefore \textit{active-low}: a low output activates an LED and a
high output deactivates it.

\begin{table}[H]
    \centering
    \caption{Relationship between the GPIO output and an active-low LED}
    \label{tab:active-low}
    \begin{tabular}{ccc}
        \toprule
        \textbf{HAL state} & \textbf{Approximate output voltage} & \textbf{LED state} \\
        \midrule
        \texttt{GPIO\_PIN\_RESET} & 0 V & On \\
        \texttt{GPIO\_PIN\_SET} & 3.3 V & Off \\
        \bottomrule
    \end{tabular}
\end{table}

% ============================================================
\subsection{Exercise 1: Alternating Two LEDs}

\subsubsection{Objective and GPIO configuration}

Exercise 1 controls a red LED connected to PA5 and a yellow LED connected to
PA6. The two LEDs must alternate every two seconds. When the red LED is on, the
yellow LED is off, and after two seconds their states are exchanged.

\begin{table}[H]
    \centering
    \caption{GPIO configuration for Exercise 1}
    \label{tab:exercise1-gpio}
    \begin{tabular}{ccc}
        \toprule
        \textbf{GPIO pin} & \textbf{CubeMX label} & \textbf{Function} \\
        \midrule
        PA5 & \texttt{LED\_RED} & Red LED output \\
        PA6 & \texttt{LED\_YELLOW} & Yellow LED output \\
        \bottomrule
    \end{tabular}
\end{table}

\subsubsection{Schematic}

Figure~\ref{fig:exercise1-schematic} is reserved for the Proteus schematic of
Exercise 1. The figure caption links to the corresponding Proteus project.

\begin{figure}[H]
    \centering
    \schematicplaceholder{figures/exercise1-schematic.png}{0.86}
    \caption{\href{\exerciseOneURL}{Exercise 1 Proteus schematic and downloadable project.}}
    \label{fig:exercise1-schematic}
\end{figure}

\subsubsection{Implementation and operation}

The program first writes a low level to PA5 to activate the red LED and a high
level to PA6 to deactivate the yellow LED. After a 2000-ms delay, the output
levels are reversed. Because these instructions are placed inside an infinite
loop, the two states repeat continuously.

\begin{stmcode}{Alternating red and yellow LEDs}
while (1)
{
    /* State 1: Red ON, Yellow OFF */
    HAL_GPIO_WritePin(LED_RED_GPIO_Port,
                      LED_RED_Pin,
                      GPIO_PIN_RESET);

    HAL_GPIO_WritePin(LED_YELLOW_GPIO_Port,
                      LED_YELLOW_Pin,
                      GPIO_PIN_SET);

    HAL_Delay(2000);

    /* State 2: Red OFF, Yellow ON */
    HAL_GPIO_WritePin(LED_RED_GPIO_Port,
                      LED_RED_Pin,
                      GPIO_PIN_SET);

    HAL_GPIO_WritePin(LED_YELLOW_GPIO_Port,
                      LED_YELLOW_Pin,
                      GPIO_PIN_RESET);

    HAL_Delay(2000);
}
\end{stmcode}

The resulting sequence is summarized in Table~\ref{tab:exercise1-sequence}.

\begin{table}[H]
    \centering
    \caption{State sequence for Exercise 1}
    \label{tab:exercise1-sequence}
    \begin{tabular}{cccc}
        \toprule
        \textbf{State} & \textbf{Red LED} & \textbf{Yellow LED} & \textbf{Duration} \\
        \midrule
        1 & On & Off & 2 seconds \\
        2 & Off & On & 2 seconds \\
        \bottomrule
    \end{tabular}
\end{table}

% ============================================================
\subsection{Exercise 2: Three-LED Traffic Light}

\subsubsection{Objective and GPIO configuration}

Exercise 2 extends the previous circuit by adding a green LED to PA7. The three
LEDs simulate a simple traffic light. The red state lasts five seconds, the
yellow state lasts two seconds, and the green state lasts three seconds.

\begin{table}[H]
    \centering
    \caption{GPIO configuration for Exercise 2}
    \label{tab:exercise2-gpio}
    \begin{tabular}{ccc}
        \toprule
        \textbf{GPIO pin} & \textbf{CubeMX label} & \textbf{Function} \\
        \midrule
        PA5 & \texttt{LED\_RED} & Red LED output \\
        PA6 & \texttt{LED\_YELLOW} & Yellow LED output \\
        PA7 & \texttt{LED\_GREEN} & Green LED output \\
        \bottomrule
    \end{tabular}
\end{table}

\subsubsection{Schematic}

\begin{figure}[H]
    \centering
    \schematicplaceholder{figures/exercise2-schematic.png}{0.86}
    \caption{\href{\exerciseTwoURL}{Exercise 2 Proteus schematic and downloadable project.}}
    \label{fig:exercise2-schematic}
\end{figure}

\subsubsection{Implementation and operation}

Only one LED is activated in each traffic-light state. Each state is written
explicitly so that the required output levels can be inspected easily. This
approach also guarantees that the two unused LEDs are switched off before the
program waits for the specified duration.

\begin{stmcode}{Three-state traffic-light controller}
while (1)
{
    /* RED for 5 seconds */
    HAL_GPIO_WritePin(GPIOA, LED_RED_Pin,
                      GPIO_PIN_RESET);
    HAL_GPIO_WritePin(GPIOA, LED_YELLOW_Pin,
                      GPIO_PIN_SET);
    HAL_GPIO_WritePin(GPIOA, LED_GREEN_Pin,
                      GPIO_PIN_SET);

    HAL_Delay(5000);

    /* YELLOW for 2 seconds */
    HAL_GPIO_WritePin(GPIOA, LED_RED_Pin,
                      GPIO_PIN_SET);
    HAL_GPIO_WritePin(GPIOA, LED_YELLOW_Pin,
                      GPIO_PIN_RESET);
    HAL_GPIO_WritePin(GPIOA, LED_GREEN_Pin,
                      GPIO_PIN_SET);

    HAL_Delay(2000);

    /* GREEN for 3 seconds */
    HAL_GPIO_WritePin(GPIOA, LED_RED_Pin,
                      GPIO_PIN_SET);
    HAL_GPIO_WritePin(GPIOA, LED_YELLOW_Pin,
                      GPIO_PIN_SET);
    HAL_GPIO_WritePin(GPIOA, LED_GREEN_Pin,
                      GPIO_PIN_RESET);

    HAL_Delay(3000);
}
\end{stmcode}

One complete traffic-light cycle lasts ten seconds, as shown in
Table~\ref{tab:exercise2-sequence}.

\begin{table}[H]
    \centering
    \caption{State sequence for Exercise 2}
    \label{tab:exercise2-sequence}
    \begin{tabular}{ccccc}
        \toprule
        \textbf{State} & \textbf{Red} & \textbf{Yellow} & \textbf{Green} & \textbf{Duration} \\
        \midrule
        1 & On & Off & Off & 5 seconds \\
        2 & Off & On & Off & 2 seconds \\
        3 & Off & Off & On & 3 seconds \\
        \bottomrule
    \end{tabular}
\end{table}

% ============================================================
\subsection{Exercise 3: Four-Way Traffic Light}

\subsubsection{Objective and control strategy}

Exercise 3 expands the circuit into a four-way traffic light containing twelve
physical LEDs. Each direction has one red, one yellow, and one green LED. In
the implemented two-phase intersection, opposite directions are synchronized:
North and South always display the same state, while East and West always
display the same state.

This design uses six GPIO control signals. Each output controls two physical
LEDs of the same color in opposite directions. Therefore, the system still
contains all twelve required LEDs:

\[
    6\ \text{GPIO control signals}
    \times 2\ \text{physical LEDs per signal}
    = 12\ \text{physical LEDs}.
\]

The North-South and East-West groups must never display green simultaneously
because their traffic paths intersect. While one group is green or yellow, the
perpendicular group remains red.

\begin{table}[H]
    \centering
    \caption{GPIO configuration for Exercise 3}
    \label{tab:exercise3-gpio}
    \begin{tabular}{ccl}
        \toprule
        \textbf{GPIO pin} & \textbf{CubeMX label} & \textbf{Physical LEDs controlled} \\
        \midrule
        PA5  & \texttt{NS\_RED}    & North red and South red \\
        PA6  & \texttt{NS\_YELLOW} & North yellow and South yellow \\
        PA7  & \texttt{NS\_GREEN}  & North green and South green \\
        PA8  & \texttt{EW\_RED}    & East red and West red \\
        PA9  & \texttt{EW\_YELLOW} & East yellow and West yellow \\
        PA10 & \texttt{EW\_GREEN}  & East green and West green \\
        \bottomrule
    \end{tabular}
\end{table}

Each physical LED should use an individual current-limiting resistor even when
two opposite LEDs share the same GPIO control signal.

\subsubsection{Schematic}

\begin{figure}[H]
    \centering
    \schematicplaceholder{figures/exercise3-schematic.png}{0.90}
    \caption{\href{\exerciseThreeURL}{Exercise 3 Proteus schematic and downloadable project.}}
    \label{fig:exercise3-schematic}
\end{figure}

\subsubsection{Traffic-light sequence}

The four-way intersection uses four states. Each direction receives a
three-second green state and a two-second yellow state. Its red state lasts
five seconds, which is equal to the other group's green and yellow durations.

\begin{table}[H]
    \centering
    \caption{State sequence for the four-way traffic light}
    \label{tab:exercise3-sequence}
    \begin{tabular}{cccc}
        \toprule
        \textbf{State} & \textbf{North-South} & \textbf{East-West} & \textbf{Duration} \\
        \midrule
        1 & Green & Red & 3 seconds \\
        2 & Yellow & Red & 2 seconds \\
        3 & Red & Green & 3 seconds \\
        4 & Red & Yellow & 2 seconds \\
        \bottomrule
    \end{tabular}
\end{table}

\subsubsection{Implementation and operation}

The program writes every output explicitly to keep the implementation easy to
read and debug. A selected color is activated with
\texttt{GPIO\_PIN\_RESET}; the other colors are deactivated with
\texttt{GPIO\_PIN\_SET}.

\begin{stmcode}{Four-way traffic-light controller}
while (1)
{
    /* State 1: North-South GREEN, East-West RED */
    HAL_GPIO_WritePin(GPIOA, NS_RED_Pin,    GPIO_PIN_SET);
    HAL_GPIO_WritePin(GPIOA, NS_YELLOW_Pin, GPIO_PIN_SET);
    HAL_GPIO_WritePin(GPIOA, NS_GREEN_Pin,  GPIO_PIN_RESET);

    HAL_GPIO_WritePin(GPIOA, EW_RED_Pin,    GPIO_PIN_RESET);
    HAL_GPIO_WritePin(GPIOA, EW_YELLOW_Pin, GPIO_PIN_SET);
    HAL_GPIO_WritePin(GPIOA, EW_GREEN_Pin,  GPIO_PIN_SET);

    HAL_Delay(3000);

    /* State 2: North-South YELLOW, East-West RED */
    HAL_GPIO_WritePin(GPIOA, NS_RED_Pin,    GPIO_PIN_SET);
    HAL_GPIO_WritePin(GPIOA, NS_YELLOW_Pin, GPIO_PIN_RESET);
    HAL_GPIO_WritePin(GPIOA, NS_GREEN_Pin,  GPIO_PIN_SET);

    HAL_GPIO_WritePin(GPIOA, EW_RED_Pin,    GPIO_PIN_RESET);
    HAL_GPIO_WritePin(GPIOA, EW_YELLOW_Pin, GPIO_PIN_SET);
    HAL_GPIO_WritePin(GPIOA, EW_GREEN_Pin,  GPIO_PIN_SET);

    HAL_Delay(2000);

    /* State 3: North-South RED, East-West GREEN */
    HAL_GPIO_WritePin(GPIOA, NS_RED_Pin,    GPIO_PIN_RESET);
    HAL_GPIO_WritePin(GPIOA, NS_YELLOW_Pin, GPIO_PIN_SET);
    HAL_GPIO_WritePin(GPIOA, NS_GREEN_Pin,  GPIO_PIN_SET);

    HAL_GPIO_WritePin(GPIOA, EW_RED_Pin,    GPIO_PIN_SET);
    HAL_GPIO_WritePin(GPIOA, EW_YELLOW_Pin, GPIO_PIN_SET);
    HAL_GPIO_WritePin(GPIOA, EW_GREEN_Pin,  GPIO_PIN_RESET);

    HAL_Delay(3000);

    /* State 4: North-South RED, East-West YELLOW */
    HAL_GPIO_WritePin(GPIOA, NS_RED_Pin,    GPIO_PIN_RESET);
    HAL_GPIO_WritePin(GPIOA, NS_YELLOW_Pin, GPIO_PIN_SET);
    HAL_GPIO_WritePin(GPIOA, NS_GREEN_Pin,  GPIO_PIN_SET);

    HAL_GPIO_WritePin(GPIOA, EW_RED_Pin,    GPIO_PIN_SET);
    HAL_GPIO_WritePin(GPIOA, EW_YELLOW_Pin, GPIO_PIN_RESET);
    HAL_GPIO_WritePin(GPIOA, EW_GREEN_Pin,  GPIO_PIN_SET);

    HAL_Delay(2000);
}
\end{stmcode}

After the fourth state, the infinite loop returns to the first state. This
creates a continuous ten-second intersection cycle without allowing the two
perpendicular traffic groups to receive green signals at the same time.

% ============================================================
\section{Conclusion}

The three exercises demonstrate progressive GPIO control using the STM32 HAL
library. Exercise 1 introduces active-low LED switching and delays. Exercise 2
extends the design into a three-state traffic light. Exercise 3 coordinates
twelve physical LEDs as a four-way intersection while using six synchronized
GPIO control signals. All three implementations use explicit output states,
which makes the programs straightforward to verify and extend in subsequent
laboratory exercises.



\end{document}
